\documentclass[11pt]{article}
\usepackage[margin=1in]{geometry}
\usepackage{amsmath,amssymb}
\usepackage{booktabs}
\usepackage{tabularx}
\usepackage{enumitem}
\usepackage[hidelinks]{hyperref}
\setlist{itemsep=2pt,topsep=4pt}
\renewcommand{\arraystretch}{1.25}

\title{Implementation Plan: Switches or Dials?\\
\large How Language Model Traits Respond to Activation Steering}
\author{Team ULP}
\date{19 November 2026}

\begin{document}
\sloppy
\maketitle

\section{Summary}

The proposal asks whether activation steering acts like a \emph{dial} (every
prompt responds gradually) or a \emph{switch} (each prompt flips at its own
threshold). We measure this with the switchiness
$S = \sigma_t^2 / (\sigma_t^2 + w^2)$ from a hierarchical probit model, across
ten traits and the Qwen3 size ladder.

Nothing in this project is trained. Steering vectors are differences of
averaged activations, and everything else is generation and scoring. We still
need GPUs because of volume (96k replies per model per vector set, each scored
by a judge) and memory (the 24B judge needs about 48\,GB).

The public repository \texttt{safety-research/persona\_vectors} is the upstream
recipe from Chen et al.\ (2025), not a pipeline for this study.
Table~\ref{tab:gap} lists what it provides and what we must build. We reuse its
trait files, its trait-generation prompt, and its difference-of-means logic
(Apache-2.0, credited), and write the rest in a new repository.

\textbf{If anyone already has pipeline code, say so} and we will build on it
instead of rewriting extraction, generation, and judging.

\begin{table}[h]
\centering
\small
\begin{tabularx}{\textwidth}{@{}XX@{}}
\toprule
Upstream has & This study needs \\
\midrule
7 traits, one prompt set each &
10 traits; 3 independent extraction sets and 120 held-out evaluation prompts per trait \\
OpenAI API judge, 0--100 scale &
Local Mistral-Small-24B judge, expected rating over digits 0--9; Qwen3-14B re-scores 10\% \\
Steering coefficient in raw units &
Strength $a$ in units of the median residual norm at the middle layer \\
Hook with three position modes (all, prompt, response) &
Extra modes for Q3: first $k$ generated tokens, tokens after $k$ \\
One coefficient per run &
Resumable grid: 10 traits $\times$ 120 prompts $\times$ 10 strengths $\times$ 8 samples \\
No statistics &
Hierarchical fit, bootstrap intervals, permutation test, greedy check \\
\bottomrule
\end{tabularx}
\caption{Gap between the upstream repository and this study.}
\label{tab:gap}
\end{table}

\section{Repository layout}

\begin{verbatim}
configs/     traits.yaml, models.yaml, grid.yaml
data/traits/<trait>/   extract_A.json, extract_B.json, extract_C.json, eval.json
steer/
  models.py    load model, chat template, middle layer, median residual norm
  hooks.py     PositionSteerer (all | prompt | response | first_k | after_k)
  extract.py   pos/neg replies -> judge filter -> difference of means
  generate.py  resumable grid runner
  judge.py     local judge, expected rating over digit probabilities
analysis/
  hierfit.py   hierarchical probit, S, prompt bootstrap
  greedy.py    greedy flip points -> check on sigma_t
  q2_groups.py decision-vs-style permutation test
  q3_commit.py switching probability vs k
  q4_scale.py  mu, S, coherence window vs model size; side-effect regimes
  figures.py
slurm/       job templates for extract / generate / judge
tests/
\end{verbatim}

\section{Steps}

\subsection{Setup}
Git repository, pinned Python environment, and the five overlapping upstream
trait files vendored as extraction set A.

\subsection{Trait data}
Trait groups are fixed before any measurement and committed to
\texttt{configs/traits.yaml}:
\begin{itemize}
  \item \textbf{Decision:} refusing, sycophantic (agrees with a false claim),
        hallucinating (fabricates), deceptive, clarifying-question.
  \item \textbf{Style:} humorous, optimistic, poetic, neurotic, impolite.
\end{itemize}
For each trait we write three extraction sets (five positive/negative
instruction pairs and 20 questions each, in the upstream format) and one
120-prompt held-out evaluation set, with a check that no evaluation prompt
appears in any extraction set. Judge rubrics are rewritten to a 0--9 scale, one
per trait, plus a coherence rubric. \textbf{Everyone should skim the prompts
for their traits.}

\subsection{Models and steering}
\begin{itemize}
  \item Qwen3 is loaded with thinking disabled, so replies are not preceded by
        a reasoning block.
  \item The steered layer is the middle layer. The median residual norm at that
        layer is measured once per model on a fixed neutral prompt set.
  \item Steering adds $a \times \text{median norm} \times \hat{v}$ to the block
        output, where $\hat{v}$ is the unit steering vector.
  \item The hook tracks prompt length and generation step, so it can steer
        prompt positions only, the first $k$ generated tokens only, or
        generated tokens after $k$.
\end{itemize}

\subsection{Extraction, generation, judging}
\begin{itemize}
  \item \textbf{Extract.} For each (model, trait, set): generate under positive
        and negative system prompts, keep pairs the judge scores as expressing
        and not expressing the trait and as coherent, average response-token
        activations, and take the difference. The first sanity result is the
        cosine similarity between sets A, B, and C.
  \item \textbf{Generate.} One job per (model, trait, vector set). Temperature
        1, 8 samples per (prompt, strength), plus one greedy sample; replies
        capped at about 200 tokens. Output is append-only, so a job killed by
        the time limit resumes where it stopped.
  \item \textbf{Judge.} Mistral-Small-24B reads each reply and its next-token
        distribution over the digits 0--9 gives the score (probability-weighted
        mean). A reply expresses the trait when its score is at least 5.
        Qwen3-14B re-scores a fixed 10\% subset; traits with Spearman
        correlation below 0.8 between judges are reported but excluded from Q2.
\end{itemize}

\subsection{Pilot on Qwen3-4B}
A coarse sweep (20 prompts, 3 traits) chooses the shared 10-point log-spaced
strength grid, plus an $a = 0$ baseline, so that it spans ``no effect'' to
``past the coherence break''. The pilot also measures throughput to confirm the
GPU budget before the full grid is launched.

\subsection{Q1: switchiness}
For prompt $i$ and strength $a$, $k$ of 8 replies express the trait. We fit
\[
  p_i(a) = \Phi\!\left(\frac{\log a - \log t_i}{w}\right),
  \qquad \log t_i \sim \mathcal{N}(\mu, \sigma_t^2),
\]
by marginal maximum likelihood, integrating numerically over each prompt's
threshold, and report $S = \sigma_t^2/(\sigma_t^2 + w^2)$ with intervals from
resampling prompts. As an independent check, the per-prompt flip strength under
greedy decoding gives a second estimate of $\sigma_t$. Qwen3-4B is run with
three vector sets; $S$ is reported per set and pooled.

\subsection{Q2: decision versus style}
An exact permutation test over trait labels on the difference in mean $S$
between the two groups (252 possible 5/5 splits).

\subsection{Q3: where the switch is decided}
Qwen3-4B, four traits (two per group), strength fixed near each trait's $\mu$.
Interventions: prompt positions only; first $k$ generated tokens; tokens after
an unsteered prefix of length $k$; $k \in \{1, 2, 4, 8, 16, 32, \text{all}\}$.
The output is switching probability against $k$ for each intervention, with
bootstrap intervals.

\subsection{Q4: scale and regimes}
\begin{itemize}
  \item \textbf{Ladder.} Qwen3 at 0.6B, 1.7B, 4B, 8B, and 14B. For each model
        we estimate $\mu$, $S$, and the coherence-break strength. The usable
        window, $\log(\text{break}) - \mu$, is regressed on log parameter
        count.
  \item \textbf{Side effects.} A reduced matrix: steer each trait vector and
        three norm-matched random directions at one strength below and one
        above the coherence break, and score 40 neutral prompts with all ten
        rubrics. The two-regime account predicts vector-specific side effects
        below the break and a shared profile above it, random directions
        included. Estimates are interpreted only where test--retest across
        vector sets has $r > 0.7$.
\end{itemize}

\subsection{Gemma 3 replication (stretch)}
4B and 12B first, then 1B and 27B if budget remains; Q1 and the window only.

\subsection{Figures}
Per-prompt curves against the population curve, $S$ per trait with intervals,
Q3 curves against $k$, window against model size, and side-effect matrices in
both regimes.

\section{Decisions to check}

These depart from, or fill gaps in, the proposal. Please object if you
disagree.

\begin{enumerate}
  \item \textbf{Baseline expression.} The proposed model forces $p \to 0$ as
        $a \to 0$, but some prompts express a trait unsteered. We add a
        per-trait baseline rate $g$ estimated with the $a = 0$ replies:
        $p = g + (1 - g)\,\Phi(\cdot)$. Without it, $w$ and $\sigma_t$ are
        biased.
  \item \textbf{Incoherent replies at high strength.} Past the coherence
        break, trait scores fall again because replies turn to noise, which
        breaks the monotone model. We fit only strengths at or below the
        trait's coherence break (where mean coherence drops below 5) and report
        the break separately.
  \item \textbf{Judging cost.} Scoring every reply for every trait is about one
        million judge calls per model and vector set. The main grid is scored
        for its own trait and coherence only; all ten rubrics are used only on
        the reduced side-effect matrix.
  \item \textbf{Vector sets on the ladder.} The full grid with three sets on
        all five models probably exceeds 250 GPU-hours once judging is counted.
        Proposed: the full grid with three sets on 4B; on the other sizes,
        three sets with 60 prompts. The pilot's throughput decides whether
        this cut is needed.
  \item \textbf{Tenth trait.} The proposal names nine traits; impolite is
        added as the fifth style trait because upstream data exists for it.
\end{enumerate}

\section{Schedule}

\begin{center}
\begin{tabular}{@{}ll@{}}
\toprule
Week & Work \\
\midrule
1 (Nov 19--25) & Pipeline, trait data, local smoke test, pilot on Engaging \\
2 (Nov 26--Dec 2) & Full 4B grid, Q1 fit, Q2 test \\
3 (Dec 3--9) & Q3 runs; ladder jobs launched \\
4 (Dec 10--16) & Q4 analysis, side-effect matrix, judge cross-check, Gemma if budget allows \\
5 & Figures and writing \\
\bottomrule
\end{tabular}
\end{center}

\section{What is needed from the team}

\begin{itemize}
  \item Engaging details: partition and account for H100s, job time limit, and
        a scratch path for model weights.
  \item A Hugging Face token on the cluster with the Gemma 3 and Mistral-Small
        licences accepted.
  \item Whether pipeline code already exists anywhere.
  \item Objections to any item in ``Decisions to check''.
\end{itemize}

\section{Verification}

\begin{itemize}
  \item \textbf{Unit tests.} The hook changes exactly the intended positions by
        exactly $a \times \text{median norm}$, in every mode; with $a = 0$ the
        output is unchanged.
  \item \textbf{Parameter recovery.} Simulate replies from known $\mu$,
        $\sigma_t$, $w$, and $g$ at the real design size and confirm the fit
        recovers $S$, for a switch-like and a dial-like setting. This also
        shows how tight the intervals on $S$ can be before any GPU time is
        spent.
  \item \textbf{Smoke test.} Qwen3-0.6B on a laptop, one trait, a few prompts
        and strengths: extract, generate, judge, and fit run end to end and
        resume after an interrupt.
  \item \textbf{Pilot.} Read 50 steered replies by hand against judge scores
        for three traits before trusting the grid, and check vector agreement
        across sets A, B, and C.
\end{itemize}

\end{document}
